% 06-ket-luan.tex — Chương 6: Kết luận và hướng phát triển
\chapter{Kết luận và hướng phát triển}
\label{ch:ketluan}

\section{Kết luận}

Báo cáo đã nghiên cứu và triển khai thực nghiệm một trợ lý ảo thông minh cho Tổng công ty
Bưu điện Việt Nam, dựa trên mô hình ngôn ngữ lớn kết hợp kỹ thuật RAG
(\rag) và fine-tuning (\ft). Các kết quả chính bao gồm:

\begin{enumerate}
  \item \textbf{Hệ thống hóa cơ sở lý thuyết:} làm rõ kiến trúc Transformer, cơ chế của
  \rag{} (chunking, embedding, hybrid retrieval, rerank), các phương pháp \ft{}
  (SFT, LoRA/QLoRA, DPO), các thước đo đánh giá và nền tảng triển khai (vLLM, FP8, MCP, A2A,
  OpenTelemetry, MLflow). Từ đó, báo cáo lập luận cho kiến trúc hybrid \rag{} + \ft.
  \item \textbf{Thiết kế kiến trúc multi-agent:} đề xuất kiến trúc gồm một agent điều phối và
  các agent chuyên biệt giao tiếp qua A2A, truy cập công cụ qua MCP, phục vụ cục bộ nhằm
  bảo đảm dữ liệu nội bộ không rời hạ tầng tổ chức.
  \item \textbf{Triển khai pipeline \rag{} lai:} cài đặt sparse retrieval BM25 và dense retrieval
  trên Qdrant, kết hợp rerank bằng cross-encoder, cùng quy trình tiền xử lý tài liệu
  (OCR, Markdown hóa, metadata hóa, structure-aware chunking).
  \item \textbf{Pipeline \ft:} đề xuất quy trình thu thập dữ liệu, ẩn PII và fine-tuning cho
  cả ba loại mô hình -- LLM (LoRA/QLoRA), embedding model (contrastive learning) và reranker
  cross-encoder -- và \emph{thí điểm} huấn luyện, đóng gói cho embedding model; kèm tiêu chí
  quyết định khi nào dùng \rag, \ft{} hoặc kết hợp. Thí điểm embedding đạt Accuracy@1 tăng từ
  0,2556 lên 0,3291 và NDCG@10 từ 0,3931 lên 0,5132 (tăng tương đối khoảng 29--31\%).
  \item \textbf{Đánh giá retrieval:} đã xây dựng quy trình đánh giá retrieval tái lập đúng
  đường đi sản xuất, kèm cơ chế phân tách lỗi giữa khâu retrieval và khâu rerank.
  \item \textbf{Hiệu năng và hiệu quả:} hạ tầng vLLM với quantization FP8 phục vụ được các
  mô hình cỡ 20--35 tỷ tham số, đủ đáp ứng nhu cầu đồng thời của một tổng công ty; giải
  pháp giúp giảm thời gian tra cứu và chuẩn hóa câu trả lời ở khâu vận hành.
\end{enumerate}

Tổng hợp lại, \rag{} là giải pháp phù hợp và khả thi cho nhu cầu tri thức động và trích
nguồn của Vietnam Post, \ft{} bổ trợ ở lớp hành vi và phong cách nghiệp vụ, và kiến trúc hybrid là
đích đến dài hạn mang lại giá trị thực tiễn cao.

\section{Hạn chế}

Bên cạnh kết quả đạt được, đề tài còn một số hạn chế:

\begin{itemize}
  \item \textbf{Chưa có đánh giá diện rộng với người dùng cuối:} kết quả hiện tại tập trung
  vào khâu retrieval và thí điểm embedding, chưa đo lường mức độ hữu ích trên người dùng thật.
  \item \textbf{Suy giảm recall ở khâu rerank:} cấu hình hiện tại làm mất một phần
  chunk đúng ở top-$k$ nhỏ, cần cải tiến.
  \item \textbf{Phần \ft{} chưa hoàn chỉnh:} embedding model đã được thí điểm và có kết quả
  định lượng, nhưng trên mô hình 0.6B và corpus TTHC độc lập (chưa phải dữ liệu sản xuất của
  Vietnam Post); LLM và reranker mới dừng ở thiết kế, chưa có bằng chứng huấn luyện thật.
  \item \textbf{Phạm vi miền hẹp:} thực nghiệm tập trung vào bưu chính/EMS; tính khái quát
  sang các miền khác chưa được kiểm chứng.
\end{itemize}

\section{Hướng phát triển}

Trên cơ sở kết quả và hạn chế, đề tài đề xuất các hướng phát triển sau:

\begin{enumerate}
  \item \textbf{Mở rộng và chuẩn hóa đánh giá:} chuẩn hóa quy trình đo retrieval trên bộ mẫu
  lớn hơn (hàng trăm tới hàng nghìn câu); bổ sung đánh giá A/B với người dùng cuối.
  \item \textbf{Cải tiến retrieval và rerank:} điều chỉnh $k_2$ và ngưỡng điểm; thử
  nghiệm kết hợp điểm rerank với điểm retrieval; mở rộng thí điểm fine-tune embedding lên mô
  hình lớn hơn (Qwen3-Embedding-8B) trên dữ liệu tiếng Việt miền bưu chính, và fine-tune
  reranker để nâng Recall@5 sau rerank.
  \item \textbf{Nâng cấp lên agentic RAG:} cho phép agent tự quyết định truy xuất nhiều vòng,
  cải viết truy vấn và gọi công cụ bổ trợ khi ngữ cảnh chưa đủ.
  \item \textbf{Khai thác GraphRAG:} xây dựng đồ thị tri thức từ tài liệu để truy xuất theo
  quan hệ giữa các thực thể (dịch vụ, biểu phí, quy định, địa bàn), hỗ trợ câu hỏi phức tạp.
  \item \textbf{Triển khai \ft{} định kỳ:} khi đã tích lũy đủ log hội thoại chất lượng, tiến
  hành SFT/LoRA và đánh giá so với \rag{} thuần; xây dựng quy trình cập nhật định kỳ.
  \item \textbf{Continual learning và cá nhân hóa:} cập nhật tri thức liên tục và học sở
  thích người dùng, đồng thời bảo đảm kiểm soát chất lượng và an toàn.
  \item \textbf{Mở rộng đa ngôn ngữ và đa kênh:} hỗ trợ thêm ngôn ngữ và tích hợp vào các
  kênh như ứng dụng di động, website, tổng đài.
  \item \textbf{Tăng cường an toàn và tuân thủ:} kiểm soát nội dung đầu ra, chống rò rỉ dữ
  liệu nhạy cảm, và bảo đảm tuân thủ quy định về dữ liệu cá nhân.
\end{enumerate}

\section*{Kết chương}
\addcontentsline{toc}{section}{Kết chương}
Chương này đã tóm tắt kết quả và đóng góp của đề tài, chỉ rõ các hạn chế cần khắc phục và
đề xuất lộ trình phát triển. Những kết quả này là cơ sở để Vietnam Post tiếp tục mở rộng hệ thống
trợ lý ảo trong tương lai.
